% Figure from MATH 451 notes, section 32. Compile with the notes preamble.
\begin{tikzpicture}
\begin{axis}[nfig, width=0.56\textwidth, height=0.34\textwidth, clip=false,
  xmin=0, xmax=2.2, ymin=0, ymax=1.45,
  xtick={0,0.65,1,2}, xticklabels={$0$,$1 - \frac\varepsilon2$,$1$,$2$},
  ytick={1}, yticklabels={$1$}, xlabel={}]
  \addplot[draw=none, fill=figB!12] coordinates {(0,0) (0.65,0) (0.65,1) (0,1)} -- cycle;
  \addplot[draw=figR, thin, fill=figR!15] coordinates {(0.65,0) (1,0) (1,1) (0.65,1)} -- cycle;
  \addplot[figB, very thick] coordinates {(0,1) (1,1)};
  \addplot[figB, very thick] coordinates {(1,0) (2,0)};
  \addplot[only marks, mark=*, mark size=1.7pt, figB] coordinates {(0,1) (1,0) (2,0)};
  \addplot[only marks, mark=*, mark size=1.7pt, mark options={fill=white, draw=figB, thick}] coordinates {(1,1)};
  \node[font=\scriptsize, figB!80!black] at (axis cs:0.33,0.5) {$L(g,P)$};
  \node[font=\scriptsize, figR, anchor=west] at (axis cs:1.08,0.55) {$U - L = \frac\varepsilon2$};
  \draw[figR!70, -stealth, thin] (axis cs:1.1,0.55) -- (axis cs:0.86,0.55);
  \node[font=\small, figB] at (axis cs:0.5,1.2) {$g$};
\end{axis}
\end{tikzpicture}
